% Junior Level 1, Week 3 -- workbook chapter (read aloud by a grown-up)
\chapterweek{3}{Bugs Are Clues}{A bug is a clue: test, look, change one thing.}

\section{What you will learn}
\begin{ican}
  \item I can test my code and look for clues.
  \item I can add a wait so I can see each step.
  \item I can fix steps in the wrong order.
\end{ican}

\section{Be a bug detective}
A \textbf{bug} is a mistake in code. Every coder finds bugs. A bug is a clue!

\begin{center}
%% alt: A circle of three steps joined by arrows. Step 1, Test, with a flag: click the green flag. Step 2, Look, with an eye: what went wrong? Step 3, Change one thing, with a wrench. An arrow goes from step 3 back to step 1.
\begin{tikzpicture}[x=1cm,y=1cm,
  stage/.style={draw=kcblack,line width=0.8pt,rounded corners=6pt,minimum width=3.6cm,minimum height=2.2cm,align=center,font=\bfseries},
  go/.style={line width=2pt,-{Triangle[length=9pt,width=12pt]},kcblue}]
  \node[stage,fill=kcgreen!15] (t) at (5,3) {{\Large\faFlag}\\1. Test};
  \node[stage,fill=kcsky!20] (l) at (10,0) {{\Large\faEye}\\2. Look};
  \node[stage,fill=kcorange!20] (c) at (0,0) {{\Large\faWrench}\\3. Change\\one thing};
  \draw[go] (t.east) to[bend left=25] (l.north);
  \draw[go] (l.west) -- (c.east);
  \draw[go] (c.north) to[bend left=25] (t.west);
\end{tikzpicture}
\end{center}

\begin{exercise}[Ready, Set, Go!]
The cat should say Ready, then Set, then Go!
\begin{enumerate}[leftmargin=1.4em,itemsep=2pt]
  \item Build this script. Use the say block with no seconds.
\end{enumerate}
\begin{center}
\begin{scratch}[scale=0.9]
\blockinit{when \greenflag clicked}
\blocklook{say \ovalnum{Ready}}
\blocklook{say \ovalnum{Set}}
\blocklook{say \ovalnum{Go!}}
\blockmove{glide \ovalnum{1} secs to \selectmenu{random position}}
\end{scratch}
\end{center}
\begin{enumerate}[leftmargin=1.4em,itemsep=2pt,start=2]
  \item Test: click the green flag.
  \item Look: did you see Ready? Did you see Set?
  \item The computer is too fast! Find the wait block in orange Control.
\end{enumerate}
\begin{center}
\begin{scratch}[scale=0.9]
\blockcontrol{wait \ovalnum{1} seconds}
\end{scratch}
\end{center}
\begin{enumerate}[leftmargin=1.4em,itemsep=2pt,start=5]
  \item Change one thing: snap a wait under Ready. Test again.
  \item Is Set still too fast? Fix it the same way.
\end{enumerate}
\end{exercise}

\needspace{14\baselineskip}
\begin{exercise}[The Sleepy Cat]
Sam's cat should sleep, then wake up, then walk. Build it and click the flag. What goes wrong?
\begin{center}
\begin{scratch}[scale=0.9]
\blockinit{when \greenflag clicked}
\blocklook{say \ovalnum{Good morning!} for \ovalnum{2} seconds}
\blockmove{move \ovalnum{50} steps}
\blocklook{think \ovalnum{Zzz} for \ovalnum{2} seconds}
\end{scratch}
\end{center}
\begin{enumerate}[leftmargin=1.4em,itemsep=2pt]
  \item Circle the block in the wrong place.
  \item Fix it on your Mac. Test again.
\end{enumerate}
\end{exercise}

\begin{tip}[Kind words]
When you find a bug in a friend's code, say: \enquote{I think I found a clue!} When a friend finds yours, say: \enquote{Thanks!}
\end{tip}

\begin{vocab}
  \vocabword{bug}{A mistake in code. It is a clue!}
  \vocabword{debug}{To find a bug and fix it.}
  \vocabword{test}{Click the green flag and watch what happens.}
\end{vocab}

\needspace{12\baselineskip}
\begin{reflection}
Draw a bug detective. How do you feel when you find a bug? Circle one.

\begin{center}
\Huge\faSmile[regular]\qquad\faMeh[regular]\qquad\faFrown[regular]
\end{center}
\answerspace{3.5cm}
\end{reflection}
